\documentclass[10pt,a4paper]{article}

\usepackage[utf8]{inputenc}

\usepackage{indentfirst}    % Красная строка на первом абзаце

% \usepackage{subcaption}
\usepackage[caption = false]{subfig}
\usepackage{hyperref}

% set font default
% \renewcommand*\familydefault{\sfdefault} 	
\usepackage[T2A,T1]{fontenc}
\usepackage[english,russian]{babel}	% локализация и переносы
\usepackage{libertine}


% more font size definitions
\usepackage{moresize}

%define page styles using geometry
\usepackage[a4paper]{geometry}		

%use customized header
\usepackage{fancyhdr}				
% \pagestyle{fancy}

%----------------------------------------------------------------------------------------
%	TABLE /ARRAY DEFINITIONS
%---------------------------------------------------------------------------------------- 

%for layouting tables
\usepackage{multicol}			
\usepackage{multirow}

\usepackage{enumitem}  % Для улучшенной настройки списка

%extended aligning of tabular cells
\usepackage{array}

\newcolumntype{x}[1]{%
>{\raggedleft\hspace{0pt}}p{#1}}%

%----------------------------------------------------------------------------------------
%	GRAPHICS DEFINITIONS
%---------------------------------------------------------------------------------------- 

%for header image
% \usepackage{graphicx}

%for floating figures
\usepackage{wrapfig}
\usepackage{float}
%\floatstyle{boxed} 
%\restylefloat{figure}

%for drawing graphics		
\usepackage{tikz}				
\usetikzlibrary{shapes, backgrounds,mindmap, trees}


%----------------------------------------------------------------------------------------
%	Color DEFINITIONS
%---------------------------------------------------------------------------------------- 

\usepackage{color}

%accent color
\definecolor{sectcol}{RGB}{255,150,0}

%dark background color
\definecolor{bgcol}{RGB}{110,110,110}

%light background / accent color
\definecolor{softcol}{RGB}{225,225,225}

\graphicspath{{img/}}

% for example, change the margins to 2 inches all round
\geometry{top=1.75cm, bottom=-.6cm, left=1.5cm, right=1.5cm} 	

\setlength{\parindent}{1cm}  % Устанавливает стандартный отступ для всех абзацев

\input{common/commands.tex}

\begin{document}
%use our custom fancy header definitions
% \pagestyle{fancy}


\section{Установка}

Ещё не известно как.

\section{Использование}

\subsection{Общий вид}

\begin{figure}[H]
    \centering
    \subfloat{\includegraphics[clip,width=0.8\linewidth]{001-Begin.png}}
    \caption{Внешний вид окна приложения}
    \label{fig:init}
\end{figure}

Окно программы на рисунке \ref{fig:init}.

\textbf{\textit{Раздел 1}} отвечает за выбор порта подключения. Выпадающий список содержит все доступные последовательные порты и пункт "Refresh"
для обновления списка. Подключаемые порты программа не видит "налету". требуется нажатие на выпадающий пунки "Refresh".
Пока подключения нет, остальные поля не активны.

\textbf{\textit{Раздел 2}} отвечает за настройки, выбор режимов работы. Переключать вкладки можно, однако поля будут активны, если подключен контроллер.

\textbf{\textit{Раздел 3, Раздел 4}} графики температуры, напряжени и тока, а также из текущие значения.

\subsection{Управление настройками}

Поля ввода имеют как минимум одну кнопку \textit{Get} или две кнопки \textit{Get}/\textit{Set} для получения или записи значения в контроллер.

Практически все значения хранятся в ОЗУ и сбрасываются на начальные при включение. Если требуется сохранить значение в энергонезависимой памяти,
то необходимо нажать кнопку \textit{Common\arrow Save Settings }

Список полей, которые не сохраняются в энергонезависимой памяти:
\begin{itemize}
    \item \textit{Common\arrow Mode} - режим работы контроллера
    \item \textit{Current Loop\arrow Current} - выходной ток в режиме Стабилизатор тока
    \item \textit{Temperature Loop\arrow T set} - температура стабилизации в режиме Стабилизатор температуры
\end{itemize}


\subsection{Вкладка Common}


\end{document}
